\section{Add Two Numbers Represented by Linked Lists}
\label{sec:ll-add-two-numbers}

\problemstatement{Two non-negative integers are stored as linked lists with
the \emph{least} significant digit first. Add them and return the sum as a
linked list in the same order.}

\begin{patternbox}
Because the digits are stored least-significant first, addition proceeds
head-to-tail with a running \textbf{carry} -- a straightforward simultaneous
walk of both lists, building the result node by node, exactly like grade-
school addition.
\end{patternbox}

\subsection*{Walk both lists with a carry}

Traverse both lists together. At each step add the two current digits (or
$0$ if a list is exhausted) plus the carry; the new digit is the sum modulo
$10$ and the new carry is the sum divided by $10$. Append the digit to the
result via a dummy-headed builder. Continue until both lists end and the
carry is zero -- a final carry produces one more leading node.

\begin{ideabox}[Least-Significant-First Makes Addition a Forward Walk]
Storing digits reversed means the carry flows in the same direction as
traversal, so no reversal or recursion is needed -- just a single forward
pass summing aligned digits with a carry, the most natural form of the
problem.
\end{ideabox}

\subsection*{Algorithm}

\begin{enumerate}
  \item Dummy-headed result; \texttt{carry = 0}.
  \item While either list has nodes or \texttt{carry} is nonzero: sum the
        available digits and carry; append \texttt{sum \% 10}; set
        \texttt{carry = sum / 10}.
  \item Return \texttt{dummy.next}.
\end{enumerate}

\subsection*{Dry run}

\dryrun{$2\to4\to3$ ($342$) and $5\to6\to4$ ($465$).}

\begin{itemize}
  \item $2{+}5{=}7$; $4{+}6{=}10\to0$ carry $1$; $3{+}4{+}1{=}8$.
  \item Result $7\to0\to8$ ($807$).
\end{itemize}

\subsection*{C++ implementation}

\begin{lstlisting}
#include <bits/stdc++.h>
using namespace std;

struct ListNode {
    int val; ListNode* next;
    ListNode(int v) : val(v), next(nullptr) {}
};

ListNode* build(const vector<int>& a) {
    ListNode dummy(0), *t = &dummy;
    for (int x : a) { t->next = new ListNode(x); t = t->next; }
    return dummy.next;
}

ListNode* addTwoNumbers(ListNode* a, ListNode* b) {
    ListNode dummy(0), *t = &dummy;
    int carry = 0;
    while (a || b || carry) {
        int sum = carry;
        if (a) { sum += a->val; a = a->next; }
        if (b) { sum += b->val; b = b->next; }
        carry = sum / 10;
        t->next = new ListNode(sum % 10);
        t = t->next;
    }
    return dummy.next;
}

void print(ListNode* h) { for (; h; h = h->next) cout << h->val << " "; cout << "\n"; }

int main() {
    print(addTwoNumbers(build({2, 4, 3}), build({5, 6, 4})));   // 7 0 8
    return 0;
}
\end{lstlisting}

\begin{complexitybox}
\textbf{Time Complexity.} $O(\max(a,b))$ -- one pass over the longer list.
\textbf{Space Complexity.} $O(\max(a,b))$ for the result list.
\end{complexitybox}

\begin{takeawaybox}
Least-significant-first storage turns list addition into a single forward
carry walk -- the cleanest arithmetic setup. Contrast with Add One's most-
significant-first order, which needs recursion or reversal to reach the
low digit first.
\end{takeawaybox}
